\chapter*{Literaturhinweise}
\addcontentsline{toc}{chapter}{Literaturhinweise}

Die ersten beiden Bücher sind grundlegende und gut lesbare Einführungen in das Thema:

E. Kahle, D. Lohse: \emph{Grundkurs Finanzmathematik}. Oldenbourg, 4.\ Auflage, 1998

B. Luderer: \emph{Starthilfe Finanzmathematik}. Teubner, 2.\ Auflage, 2003

\medskip
Das folgende Buch behandelt das ganze übliche Spektrum der in Grundvorlesungen gelehrten Wirtschaftsmathematik -- einschließlich der Finanzmathematik:

E. Salomon, W. Poguntke: \emph{Wirtschaftsmathematik}. Fortis Verlag, 2.\ Auflage, 2003

\medskip
Die nächste Empfehlung ist mathematisch recht anspruchsvoll und geht über den in dieser Lerneinheit dargestellten Grundstoff hinaus, insofern auch Aktien und Optionen behandelt werden. Das Buch behandelt auch didaktische Fragen der Vermittlung der dargestellten Inhalte:

P. Daume: \emph{Finanzmathematik im Unterricht}. Teubner, 2009
